% Generated draft: requires mathematical and source review.
\documentclass[11pt]{article}
\usepackage[margin=25mm]{geometry}
\usepackage{fontspec}
\setmainfont{Latin Modern Roman}
\usepackage{amsmath,amssymb,mathtools,unicode-math}
\setmathfont{Latin Modern Math}
\usepackage{enumitem,needspace,xurl,hyperref}
\setlength{\parindent}{0pt}
\setlength{\parskip}{5pt}
\setlength{\emergencystretch}{3em}
\newcommand{\N}{\mathbb{N}}
\newcommand{\Z}{\mathbb{Z}}
\newcommand{\Q}{\mathbb{Q}}
\newcommand{\R}{\mathbb{R}}
\newcommand{\C}{\mathbb{C}}
\newcommand{\F}{\mathbb{F}}
\newcommand{\A}{\mathbb{A}}
\newcommand{\PP}{\mathbb{P}}
\newcommand{\E}{\mathbb{E}}
\newcommand{\eps}{\varepsilon}
\newcommand{\dd}{\,\mathrm{d}}
\begin{document}
\section*{UnsolvedMath ID: 800707}
\textbf{Source problem code:} catalogue entry
\subsection{Definitions and mathematical statement}
Let \(F\) be a base field and let \(K/F\) range over field extensions. A central simple \(K\)-algebra is a finite-dimensional associative unital \(K\)-algebra whose center is \(K\) and which has no nonzero proper two-sided ideals. Its degree is \(\deg(A)=\sqrt{\dim_K A}\). Write \(\mathrm{CSA}_n(K)\) for the set of isomorphism classes of central simple \(K\)-algebras of degree \(n\). This is naturally identified with \(H^1(K,\operatorname{PGL}_n)\).

For a functor \(\mathcal T\) from field extensions of \(F\) to sets and an element \(x\in\mathcal T(K)\), define
\[
\operatorname{ed}_F(x)=\inf\left\{\operatorname{trdeg}_F K_0: F\subseteq K_0\subseteq K,\ x\in\operatorname{im}\bigl(\mathcal T(K_0)\longrightarrow\mathcal T(K)\bigr)\right\},
\qquad
\operatorname{ed}_F(\mathcal T)=\sup_{K/F}\sup_{x\in\mathcal T(K)}\operatorname{ed}_F(x).
\]
For a prime \(p\), essential \(p\)-dimension permits a finite extension of degree prime to \(p\) before descent:
\[
\operatorname{ed}_{F,p}(x)=\inf_{\substack{K'/K\text{ finite}\\p\nmid[K':K]}}\operatorname{ed}_F(x_{K'}),
\qquad
\operatorname{ed}_{F,p}(\mathcal T)=\sup_{K/F}\sup_{x\in\mathcal T(K)}\operatorname{ed}_{F,p}(x).
\]
For a finite group \(G\), interpreted as a constant group scheme over \(F\), write \(\operatorname{ed}_F(G)=\operatorname{ed}_F(K\mapsto H^1(K,G))\).

\textbf{Qualification concerning crossed products.} The supplied record does not define \(\operatorname{CPAlg}_n\), \(\mathrm{Alg}_G\), or the hypotheses on \(F\) and \(G\). One base-change-stable interpretation is the following. A \(G\)-Galois \(K\)-algebra is a finite étale \(K\)-algebra \(L\) of dimension \(|G|\), equipped with a \(G\)-action such that the induced action on its geometric points is simply transitive. Under an unmarked convention, \(\mathrm{Alg}_G(K)\) consists of isomorphism classes of central simple algebras of degree \(|G|\) containing such an \(L\) as a maximal étale subalgebra; the choice of \(L\), its embedding, and its action are not retained. Then \(\operatorname{CPAlg}_n(K)\) consists of the classes belonging to \(\mathrm{Alg}_G(K)\) for some finite group \(G\) of order \(n\).

These conventions are explicit provisional interpretations, not definitions verified from the cited source. A definition retaining the embedded Galois algebra or its group action gives a different moduli problem. Requiring a Galois field extension instead of a Galois étale algebra also changes the formulation and requires care under scalar extension.

\textbf{Problem, conditional on confirming the intended conventions and hypotheses.} For positive integers \(n\) and finite groups \(G\), determine bounds for
\[
\operatorname{ed}_F(\operatorname{CPAlg}_n)
\quad\text{and}\quad
\operatorname{ed}_F(\mathrm{Alg}_G),
\]
and determine
\[
\operatorname{ed}_F(\mathrm{Alg}_G)-\operatorname{ed}_F(G)
\]
whenever this difference is defined. Determine whether the implication
\[
\operatorname{ed}_F(\operatorname{CPAlg}_n)
 =\operatorname{ed}_F(\operatorname{PGL}_n)
\quad\Longrightarrow\quad
\bigl(\forall K/F\bigr)\bigl(\forall[A]\in\mathrm{CSA}_n(K)\bigr)\,
[A]\in\operatorname{CPAlg}_n(K)
\]
holds. Thus the conclusion concerns every degree-\(n\) central simple algebra over every extension of the same base field \(F\).

The introductory reference to essential \(p\)-dimension supplies context but specifies neither a prime nor a precise additional \(p\)-dimension question. No such additional assertion is inferred here.
\subsection{Short English statement}
The problem asks for essential-dimension bounds for crossed-product algebras, their relation to the essential dimension of a finite group, and whether equality with the essential dimension of \(\operatorname{PGL}_n\) forces every degree-\(n\) algebra over every extension of \(F\) to be a crossed product. The precise crossed-product conventions and base-field hypotheses require confirmation; the displayed formulation is conditional on them.
\subsection{Sources}
\begin{itemize}
\item Local UnsolvedMath registry: permanent ID \(800707\), exact source problem code \texttt{catalogue entry}, title ``Open Problems on Central Simple Algebras --- Problem 17.'' The supplied record attributes the problem to Auel, Brussel, Garibaldi, and Vishne, with proposed year \(2010\).
\item Auel, Brussel, Garibaldi, and Vishne, \emph{Open Problems on Central Simple Algebras} (2010), Problem 17. Supplied source URL: \url{https://arxiv.org/abs/1006.3304}. The supplied provenance records access on \texttt{2026-10-16} and extraction method \texttt{arxiv\_latex}; that access is reported by the record and was not independently performed here.
\item The registry labels the problem ``open'' and records an update at \texttt{2026-10-16T14:00:00Z}. This is supplied dated registry information, not an independently verified assessment of the literature or present mathematical status.
\end{itemize}
\end{document}
